\section{Conclusion}
\label{sec:conclusion}

We audited what four published coding-agent studies release, against what they report. A
codebook frozen before any availability lookup was instantiated into 206 recorded evidence
requirements across 20 comparison decisions nested in those four studies, and 2.8\,GB was
retrieved. The released bytes support four bounded operations: SWE-agent task-manifest
equality, selected-subset SWE-Gym Moatless resolved/empty arithmetic, inspection of the
Agentless generation archive, and AutoCodeRover run-report and candidate arithmetic.
They do not support a claim that every study was fully recomputed. Around them, a bounded
discovery cohort of 47 rostered papers characterises the surrounding practice.

Under the outcome-specific v2 re-adjudication, all 20 audited decisions are
\textsc{unidentified}; none is \textsc{identified} or \textsc{compatibility not evaluated},
and all five condition groups across the four studies are \textsc{not identified}. The earlier
v1 result of 18 \textsc{unidentified} plus two set aside after a 13{,}200-combination upper
bound is historical and superseded. Tier-two per-class compatibility was not run, so the paper
reports no transfer or attemptability result. The v2 verdicts reflect incomplete
outcome-specific public evidence chains; they do not establish that filling any one field
would make reproduction succeed. In the surrounding cohort, 20 of the 47 rostered papers
expose an artifact of their own and 13 of those publish run outputs -- the disposition of this
cohort, not a rate in the literature.
How many of them bind those outputs to their reported numbers under the joint criterion is
not answered here; the prose-level grades we collected do not settle it.
No actual deposit was established in advance as a known-positive control under the v2 rule,
so the zero identified decisions cannot calibrate the audit's sensitivity to a real positive.

The audit also found that reportability and recoverability are independent properties. A
task manifest coded \emph{not reported} from the paper text equals the pinned public revision in the
proceedings supplement; that check establishes manifest equality only. Conversely, released
run records do not identify a reported outcome when the required public evidence chain is
incomplete.

The bounded checks expose distinctions that the source and artifact do not resolve together.
AutoCodeRover reports three repetitions of its base configuration, cites the standard $k=1$
estimator, and prints 57/66/78, annotating only the two base figures with a run count. Public run-output reports expose base-run triples of 54/57/57 and 48/47/50 and
ACR-sbfl counts of 66 and 61. Raw per-repetition contrasts span 9--14; complete-bundle
source-unbound differences are 10 and $38/3$. The source-to-public-bundle binding and complete
selection/count presentation transform are unreported. CMP-035 aggregates 78 unique tasks with
any plausible patch and 51 with any correct patch, but its task-level human semantic-label
vector is absent. These observations select no effect and do not change the v2 verdicts.

The practical requirement is to publish an outcome-specific evidence chain that identifies
and binds the release, run outputs, evaluator configuration, task universe, aggregation, and
presentation steps needed for each reported value. This audit does not estimate the effort
needed to supply that chain or claim that any one added mapping would suffice.
